\documentclass[../../main.tex]{subfiles}
\begin{document}

% 年度の大きな表示
\begin{center}
  {\fontsize{48pt}{58pt}\selectfont \textbf{2004}\normalsize\textbf{年度}}
\end{center}

\vspace{10mm}

% 本文


\vspace{15mm}

% 出題テーマと難易度の表
\noindent\textbf{出題テーマと難易度}

\vspace{5mm}

\begin{center}
  \shadowbox{%
    \begin{tabular}{|c|p{8cm}|c|}
      \hline
      \textbf{問題番号} & \textbf{テーマ} & \textbf{難易度} \\
      \hline
      第1問 & 分数関数の増減とグラフの共有点条件 & \\
      \hline
      第2問 & 定積分の周期性と絶対値付き三角関数の極限 & \\
      \hline
      第3問 & 確率（コインの選択と玉の色の反復試行） & \\
      \hline
      第4問 & 空間図形（2球の共通部分の体積の最大化） & \\
      \hline
    \end{tabular}%
  }
\end{center}

\vspace{15mm}

% 解答
\noindent\textbf{解答}

\vspace{5mm}

\begin{center}
  \shadowbox{%
    \renewcommand{\arraystretch}{1.6}
    \begin{tabular}{|c|c|p{8cm}|}
      \hline
      \textbf{問題番号} & \textbf{小問} & \textbf{答え} \\
      \hline
      \multirow{2}{*}{第1問} & (1) & $a<x<4a$で減少，$4a<x$で増加（最小値$f(4a)=\dfrac{4^4}{3^3}a$） \\
      \cline{2-3}
                             & (2) & $b>\dfrac{2\cdot4^4}{3^4}a$ \\
      \hline
      \multirow{3}{*}{第2問} & (1) & 証明問題 \\
      \cline{2-3}
                             & (2) & 証明問題 \\
      \cline{2-3}
                             & (3) & $\dfrac{1}{4}+\dfrac{1}{2\pi}$ \\
      \hline
      \multirow{3}{*}{第3問} & (1) & $A={}_nC_k\,p^n\,q^k(1-q)^{n-k}$ \\
      \cline{2-3}
                             & (2) & $B=\{pq+(1-p)r\}^n$ \\
      \cline{2-3}
                             & (3) & 赤玉$701$個 \\
      \hline
      \multirow{2}{*}{第4問} & (1) & $V(r)=\left(-r^4+\dfrac{2}{3}r^3+r^2-\dfrac{2}{3}+\dfrac{2}{3}(1-r^2)^{3/2}\right)\pi$ \\
      \cline{2-3}
                             & (2) & $r=\dfrac{\sqrt{2}}{2}$で最大値$\left(\dfrac{\sqrt{2}}{3}-\dfrac{5}{12}\right)\pi$ \\
      \hline
    \end{tabular}%
    \renewcommand{\arraystretch}{1.0}
  }
\end{center}

\vspace{15mm}

% 解答時間
\noindent\textbf{試験形態}

\vspace{5mm}

\begin{center}
  \shadowbox{%
    \begin{tabular}{p{8cm}}
      （要確認）
    \end{tabular}%
  }
\end{center}

\end{document}
